The kernel: $(A, B)$ acts trivially iff $A \otimes B = 1$ iff $A = \lambda I_{2}$, $B = \mu I_{3}$ with $\lambda\mu = 1$, which is $K$, the kernel of \claimref{standard-model.gauge-algebra-minimal-massive-composite}.

The first isomorphism: each $K$-coset contains exactly one pair with $\det A \cdot \det B = 1$. Replacing $(A, B)$ by $(\lambda A, \lambda^{-1}B)$ rescales the product determinant by $\lambda^{2}\lambda^{-3} = \lambda^{-1}$, so within a coset the product determinant takes every value of the form $\lambda^{-1}d_{0}$ exactly once as $\lambda$ runs over $U(1)$, and $\lambda = d_{0}$ is the unique choice landing it at $1$. The inclusion $S(U(2) \times U(3)) \hookrightarrow U(2) \times U(3)$ followed by the quotient by $K$ is therefore a bijective continuous homomorphism of compact groups, hence an isomorphism: a continuous bijection from a compact space onto a Hausdorff space is a homeomorphism, since it carries closed sets, which are compact, to compact sets, which are closed.

The second isomorphism: the map $(a, b, \lambda) \mapsto (\lambda^{3}a,\, \lambda^{-2}b)$ from $SU(2) \times SU(3) \times U(1)$ lands in $S(U(2) \times U(3))$, since $\det(\lambda^{3}a)\det(\lambda^{-2}b) = \lambda^{6}\lambda^{-6} = 1$. It is surjective: given $(A, B)$ with $\det A \cdot \det B = 1$, write $\det A = e^{i\alpha}$ and pick any sixth root $\lambda^{6} = e^{i\alpha}$; then $a = \lambda^{-3}A$ has $\det a = \lambda^{-6}e^{i\alpha} = 1$ and $b = \lambda^{2}B$ has $\det b = \lambda^{6}\det B = e^{i\alpha}\det B = 1$. Its kernel is $\lambda^{3}a = I_{2}$, $\lambda^{-2}b = I_{3}$ with $a \in SU(2)$, $b \in SU(3)$: taking determinants, $\lambda^{6} = 1$ and $\lambda^{-6} = 1$, so $\lambda$ is a sixth root of unity and $(a, b) = (\lambda^{-3}I_{2}, \lambda^{2}I_{3})$ is determined by it: the cyclic group $\mathbb{Z}_{6}$, embedded as $\lambda \mapsto (\lambda^{-3}I_{2}, \lambda^{2}I_{3}, \lambda)$. The induced map from the quotient by this kernel is again a bijective continuous homomorphism of compact groups, hence an isomorphism by the same fact.

Consequence (i). The block-diagonal elements of $SU(5)$ preserving $\C^{2} \oplus \C^{3}$ are the pairs $(A, B) \in U(2) \times U(3)$ with $\det A \cdot \det B = 1$, which is $S(U(2) \times U(3))$ by definition. A cover $(SU(2) \times SU(3) \times U(1))/\Gamma$ with $\Gamma \subsetneq \mathbb{Z}_{6}$ maps onto $S(U(2) \times U(3))$ through the second isomorphism with kernel $\mathbb{Z}_{6}/\Gamma \neq 1$, so it does not act faithfully on $\C^{5}$ through the block embedding.

Consequence (ii), the lattice. A representation of the covering product $SU(2) \times SU(3) \times U(1)$ descends to the quotient exactly when the embedded $\mathbb{Z}_{6}$ acts trivially. Its generator is $\gamma = (\lambda^{-3}I_{2}, \lambda^{2}I_{3}, \lambda)$ with $\lambda = e^{2\pi i/6}$, from the kernel computation above. On an irreducible representation, $\lambda^{-3}I_{2} = -I_{2}$ acts by $(-1)^{\delta}$ where $\delta$ is the weak parity ($\delta = 1$ on a doublet, $0$ on a singlet or triplet), because the centre $\{\pm I_{2}\}$ of $SU(2)$ acts on an irreducible of dimension $r$ by $(-1)^{r-1}$; $\lambda^{2}I_{3} = e^{2\pi i/3}I_{3}$ acts by $e^{2\pi i t/3}$ where $t$ is the triality (the $\mathbb{Z}_{3}$ class of the $SU(3)$ content: $t = 1$ on $3$, $-1$ on $\bar{3}$, $0$ on $1$ and $8$), because the centre of $SU(3)$ acts on an irreducible of highest weight $(p_{1}, p_{2})$ by $e^{2\pi i(p_{1} - p_{2})/3}$; and $\lambda$ acts by $\lambda^{q}$ where $q$ is the abelian weight. So $\gamma$ acts by $\exp 2\pi i\,(\delta/2 + t/3 + q/6)$, and triviality is $2t + 3\delta + q \equiv 0 \pmod 6$, which with $Y := q/6$ is the integrality of $t/3 + \delta/2 + Y$. This is the standard reading of the quotient \cite{BaezHuerta2010, Tong2017}, proved here from the explicit generator. The script realises $\lambda$ exactly in the cyclotomic field $\mathbb{Q}(\zeta_{6})$ and confirms that $\lambda$ has order exactly six and that each of the six powers of $\gamma$ lies in the kernel (that nothing else does is the determinant argument above: $\lambda^{6} = 1$ has six solutions); it also checks the lattice condition on every representation used in Section \ref{sec:content}, and that a colour singlet of hypercharge $1/3$ fails it. The coset rescaling $\lambda^{2}\lambda^{-3} = \lambda^{-1}$ is confirmed on a unit-modulus Gaussian rational.
